\section{Majorana covariance matrices}
\label{sec:majorana}
For $n\geq0$, use coordinates $(x_1,\ldots,x_n,y_1,\ldots,y_n)$ on $\R^{2n}$; for
each mode $j$ we label its two Majorana operators $x_j$ and $y_j$. All matrices are
real unless stated otherwise, and $\|v\|^2=v^{\T}v$ is the squared Euclidean norm.
For a real symmetric matrix $A$, the notation $A\succeq0$ means $v^{\T}Av\geq0$ for
every real column vector $v$. For a complex matrix $H$, $H\succeq0$ means that
$H=H^*$ and $z^*Hz\geq0$ for every complex column vector $z$, where $H^*$ is the
conjugate transpose. We regard real matrices as complex matrices when they enter
the complex inequality. The case $n=0$ is included throughout: the empty matrix
satisfies every condition below.

\motiv{The two conditions of Definition~\ref{def:admissible} are the ones a state imposes. If
$c_1,\ldots,c_{2n}$ are Hermitian operators with $c_jc_k+c_kc_j=2\delta_{jk}$ and
$\rho$ is a density operator, the numbers $\Gamma_{jk}=\tfrac{\ii}{2}\tr(\rho[c_j,c_k])$
are real and antisymmetric, and $1+\ii\Gamma\succeq0$ is exactly the statement
$\tr(\rho A^*A)\geq0$ for the operators $A=\sum_j\overline{z_j}\,c_j$. Section~\ref{sec:fock}
proves both facts (Theorems~\ref{thm:reality} and~\ref{thm:admissible});
this section derives only what follows from the matrix conditions.}

\begin{definition}[Admissible covariance]\label{def:admissible}
A real $2n\times2n$ matrix $\Gamma$ is called \emph{admissible} (a physical covariance) if
\begin{equation}\label{eq:admissibility}
 \Gamma^{\T}=-\Gamma,\qquad 1-\Gamma^{\T}\Gamma\succeq0 .
\end{equation}
\end{definition}
The matrix $1-\Gamma^{\T}\Gamma$ is symmetric for every $\Gamma$, so the second
condition is the real positive semidefiniteness of a symmetric matrix.

\motiv{\revs{A second condition singles out the covariances of pure states: for a Gaussian state,
$\Gamma^2=-1$ holds exactly when the state is pure~\cite[Sec.~2]{bravyi}
(Theorem~\ref{thm:idempotent} in Section~\ref{sec:fmps}). Here the condition is taken as the
definition.}}

\begin{definition}[Pure covariance]\label{def:pure}
A real $2n\times2n$ matrix $\Gamma$ is called \emph{pure} if
$\Gamma^{\T}=-\Gamma$ and $\Gamma^2=-1$.
\end{definition}
\revs{A pure covariance is orthogonal, since $\Gamma^{\T}\Gamma=-\Gamma^2=1$, and it is a complex
structure on $\R^{2n}$: since $(\Gamma-\ii)(\Gamma+\ii)=\Gamma^2+1=0$, $\Gamma$ is diagonalizable
over $\C$ with eigenvalues $\pm\ii$ (a matrix annihilated by a polynomial with distinct roots is
diagonalizable), and complex conjugation maps the $\ii$-eigenspace onto the $-\ii$-eigenspace
because $\Gamma$ is real, so each has dimension $n$.}

For any real matrix $A$ and real vector $v$, $v^{\T}Av=v^{\T}A^{\T}v$; hence
$v^{\T}\Gamma v=0$ when $\Gamma$ is antisymmetric, and $a^{\T}\Gamma b=-b^{\T}\Gamma a$
for all real $a,b$. For any real matrix $O$ and real symmetric $A\succeq0$, the
congruence $OAO^{\T}$ is symmetric and
$v^{\T}OAO^{\T}v=(O^{\T}v)^{\T}A(O^{\T}v)\geq0$, so $OAO^{\T}\succeq0$.

\subsection{Equivalent forms of admissibility}

\begin{theorem}[Norm form]\label{thm:norm-form}
Let $\Gamma^{\T}=-\Gamma$. Then $\Gamma$ is admissible if and only if
$\|\Gamma v\|\leq\|v\|$ for every $v\in\R^{2n}$.
\end{theorem}
\begin{proof}
For real $v$, $v^{\T}(1-\Gamma^{\T}\Gamma)v=\|v\|^2-\|\Gamma v\|^2$.
\end{proof}

\motiv{The real form is the one the
operations of Section~\ref{subsec:operations} preserve
(Theorems~\ref{thm:orthogonal}--\ref{thm:sum}); the Hermitian form is the one a state
delivers, since $z^*(1+\ii\Gamma)z=\tr(\rho A^*A)$ for $A=\sum_j\overline{z_j}c_j$
(Theorem~\ref{thm:admissible}).}

\begin{theorem}[Hermitian form]\label{thm:hermitian-form}
Let $\Gamma^{\T}=-\Gamma$. Then $\Gamma$ is admissible if and only if
$1+\ii\Gamma\succeq0$ as a complex matrix.
\end{theorem}
\begin{proof}
The matrix $1+\ii\Gamma$ is Hermitian, since
$(\ii\Gamma)^*=-\ii\Gamma^{\T}=\ii\Gamma$. Write $z=a+\ii b$ with $a,b$ real.
Then $z^*=a^{\T}-\ii b^{\T}$ and
\[
 z^*(1+\ii\Gamma)z
 =(a^{\T}-\ii b^{\T})(a+\ii b)+\ii(a^{\T}-\ii b^{\T})\Gamma(a+\ii b)
 =\|a\|^2+\|b\|^2+\ii\bigl(\ii\,a^{\T}\Gamma b-\ii\,b^{\T}\Gamma a\bigr),
\]
using $a^{\T}\Gamma a=b^{\T}\Gamma b=0$. With $a^{\T}\Gamma b=-b^{\T}\Gamma a$ this gives
\begin{equation}\label{eq:hermitian-form}
 z^*(1+\ii\Gamma)z=\|a\|^2+\|b\|^2+2\,b^{\T}\Gamma a .
\end{equation}
Suppose $\|\Gamma a\|\leq\|a\|$ for all $a$. Expanding $0\leq\|b+\Gamma a\|^2$ gives
$2\,b^{\T}\Gamma a\geq-\|b\|^2-\|\Gamma a\|^2$, so the right side of
\eqref{eq:hermitian-form} is at least $\|a\|^2-\|\Gamma a\|^2\geq0$. Conversely, if
the right side is nonnegative for all $a,b$, the choice $b=-\Gamma a$ gives
$\|a\|^2-\|\Gamma a\|^2\geq0$. Theorem~\ref{thm:norm-form} completes both directions.
\end{proof}
\begin{remark}\label{rem:bosonic}
For bosonic covariances the corresponding condition is $\gamma-\ii\Omega_n\succeq0$
with $\Omega_n$ the symplectic form on the quadratures; for fermions the identity and
the symplectic form exchange roles. The coordinates $x_j,y_j$ of these notes label the
two Majorana operators of a mode and are not quadratures.
\end{remark}

\begin{theorem}[Purity]\label{thm:purity}
Every pure covariance is admissible, satisfies $\Gamma^{\T}\Gamma=1$, and is an
isometry: $\|\Gamma v\|=\|v\|$ for all $v$.
\end{theorem}
\begin{proof}
$\Gamma^{\T}\Gamma=-\Gamma^2=1$, so $1-\Gamma^{\T}\Gamma=0\succeq0$, and
$\|\Gamma v\|^2=v^{\T}\Gamma^{\T}\Gamma v=\|v\|^2$.
\end{proof}

\subsection{Changes of basis, subsystems and direct sums}
\label{subsec:operations}
\noindent A unitary $U$ that maps each Majorana operator to a real linear combination,
$Uc_jU^*=\sum_kO_{jk}c_k$ with $O$ real, has $O$ orthogonal, because the images again
satisfy $c_jc_k+c_kc_j=2\delta_{jk}$; the state $U^*\rho U$ then has covariance
$O\Gamma O^{\T}$~\cite[Sec.~2]{bravyi} ($U$ implements $O$ in the terminology of
Section~\ref{sec:hamiltonian}). Admissibility and purity survive this congruence
(Theorem~\ref{thm:orthogonal}). Discarding modes restricts $\Gamma$ to the remaining
coordinates; admissibility survives (Theorem~\ref{thm:subsystem}), purity need not
(Remark~\ref{rem:purity-subsystem}). Preparing two systems independently gives the direct
sum of their covariances, and admissibility and purity are decided blockwise
(Theorem~\ref{thm:sum}). We build matrix product states from these three operations
and a fourth, contraction (Section~\ref{sec:contraction}).

\begin{theorem}[Orthogonal congruence]\label{thm:orthogonal}
Let $O$ be a real $2n\times2n$ matrix with $OO^{\T}=1$. If $\Gamma$ is admissible,
then $O\Gamma O^{\T}$ is admissible; if $\Gamma$ is pure, then $O\Gamma O^{\T}$ is pure.
\end{theorem}
\begin{proof}
Since $O$ is square, $OO^{\T}=1$ makes $O$ invertible with inverse $O^{\T}$, so
$O^{\T}O=1$ as well. Antisymmetry:
$(O\Gamma O^{\T})^{\T}=O\Gamma^{\T}O^{\T}=-O\Gamma O^{\T}$. Inserting $O^{\T}O=1$,
$(O\Gamma O^{\T})^{\T}(O\Gamma O^{\T})=O\Gamma^{\T}\Gamma O^{\T}$, hence
$1-(O\Gamma O^{\T})^{\T}(O\Gamma O^{\T})=O(1-\Gamma^{\T}\Gamma)O^{\T}\succeq0$
by the congruence fact above. For purity,
$(O\Gamma O^{\T})^2=O\Gamma^2O^{\T}=-OO^{\T}=-1$.
\end{proof}

To select a subsystem, we take both Majorana coordinates of each chosen mode
together. An injective map $e:\{1,\ldots,m\}\to\{1,\ldots,n\}$ lists $m$ distinct
modes; write $\bar e(x_a)=x_{e(a)}$ and $\bar e(y_a)=y_{e(a)}$ for the induced map on
coordinates.

\begin{theorem}[Mode selection]\label{thm:subsystem}
Let $e$ be injective and let $\Gamma$ be admissible. Then the $2m\times2m$ matrix
$\Gamma'_{\alpha\beta}=\Gamma_{\bar e(\alpha)\bar e(\beta)}$ is admissible. The empty
selection $m=0$ is included.
\end{theorem}
\begin{proof}
Antisymmetry passes to the restriction, since
$\Gamma'_{\beta\alpha}=\Gamma_{\bar e(\beta)\bar e(\alpha)}=-\Gamma_{\bar e(\alpha)\bar e(\beta)}$.
For $v\in\R^{2m}$ let $\tilde v\in\R^{2n}$ have $\tilde v_{\bar e(\alpha)}=v_\alpha$ and
all other entries zero; this is well defined because $\bar e$ is injective, and
$\|\tilde v\|=\|v\|$. For each $\alpha$,
$(\Gamma'v)_\alpha=\sum_\beta\Gamma_{\bar e(\alpha)\bar e(\beta)}v_\beta=(\Gamma\tilde v)_{\bar e(\alpha)}$.
A sum of squares over the coordinates $\bar e(\alpha)$ is at most the sum over all
coordinates, so
$\|\Gamma'v\|^2\leq\|\Gamma\tilde v\|^2\leq\|\tilde v\|^2=\|v\|^2$ by
Theorem~\ref{thm:norm-form}, which then also gives admissibility of $\Gamma'$.
\end{proof}
\begin{remark}\label{rem:purity-subsystem}
Purity does not pass to subsystems: the reduced covariance of an entangled pure
state is not pure.
\end{remark}

For $\Gamma_1$ of size $2n_1$ and $\Gamma_2$ of size $2n_2$, the
\emph{direct sum} $\Gamma_1\oplus\Gamma_2$ (the covariance of two uncorrelated systems) is the $2(n_1+n_2)$ matrix with
entries $(\Gamma_1)$ on the coordinates of the first $n_1$ modes, $(\Gamma_2)$ on
those of the last $n_2$ modes, and zero cross entries, in the $(x,y)$ ordering of
the composite system.

\begin{theorem}[Independent sums]\label{thm:sum}
$\Gamma_1\oplus\Gamma_2$ is admissible if and only if $\Gamma_1$ and $\Gamma_2$ are
admissible, and pure if and only if both are pure.
\end{theorem}
\begin{proof}
Conjugation by the permutation matrix that groups the coordinates of each subsystem
together carries $\Gamma_1\oplus\Gamma_2$ to $\operatorname{diag}(\Gamma_1,\Gamma_2)$.
This conjugation is a ring isomorphism of matrices that commutes with
transposition and fixes $1$; it therefore carries antisymmetry,
$1-\Gamma^{\T}\Gamma$ and $\Gamma^2$ to the corresponding block-diagonal
quantities, and it preserves $\succeq0$ by the congruence fact with $O$ the
permutation matrix or its transpose. It remains to decide each property blockwise.
Transposition and products of block-diagonal matrices are blockwise, so
$\operatorname{diag}(\Gamma_1,\Gamma_2)$ is antisymmetric, or squares to $-1$,
exactly when both blocks do. For $v=(v_1,v_2)$,
\[
 v^{\T}\operatorname{diag}(A_1,A_2)v=v_1^{\T}A_1v_1+v_2^{\T}A_2v_2,
\]
so $\operatorname{diag}(A_1,A_2)\succeq0$ when both blocks are, and conversely
taking $v_2=0$ or $v_1=0$ isolates each block. Apply this with
$A_i=1-\Gamma_i^{\T}\Gamma_i$.
\end{proof}

\subsection{The standard covariance and one mode}
Let
\begin{equation}\label{eq:standard}
 \Omega_n=\begin{pmatrix}0&1_n\\-1_n&0\end{pmatrix}
\end{equation}
in the coordinates $(x,y)$. It is antisymmetric and $\Omega_n^2=-1$, so $\Omega_n$
is pure. On Fock
space, $\Omega_n$ is the covariance of the state with every mode occupied and $-\Omega_n$
that of the vacuum (Corollary~\ref{cor:vacuum}).
\motiv{The covariance of every state with a definite occupation in each mode is
$\Omega_n$ with the sign of each mode's $2\times2$ block chosen independently
(Theorem~\ref{thm:occupation}); Proposition~\ref{prop:one-mode} shows that these signed
blocks are the only pure one-mode covariances.}

\begin{proposition}[One mode]\label{prop:one-mode}
Every real antisymmetric $2\times2$ matrix has the form
$\Gamma_\lambda=\lambda\Omega_1$ with $\lambda\in\R$. $\Gamma_\lambda$ is admissible
if and only if $|\lambda|\leq1$, and pure if and only if $|\lambda|=1$.
\end{proposition}
\begin{proof}
Antisymmetry forces zero diagonal entries and opposite off-diagonal entries, which
is the stated form with $\lambda$ the $(x_1,y_1)$ entry. Then
$\Gamma_\lambda^{\T}\Gamma_\lambda=\lambda^2\cdot1$ and
$\Gamma_\lambda^2=-\lambda^2\cdot1$. For a real number $c$, the matrix $c\cdot1$
satisfies $c\cdot1\succeq0$ exactly when $c\geq0$, since $v^{\T}(c\cdot1)v=c\|v\|^2$;
with $c=1-\lambda^2$ this gives admissibility exactly when $\lambda^2\leq1$. Purity
is $-\lambda^2\cdot1=-1$, that is $\lambda^2=1$.
\end{proof}
